\documentclass[11pt]{article}
% Standalone preamble for web-only supplementary notes.
\usepackage[T1]{fontenc}
\usepackage{amsmath,amssymb,amsthm}
\usepackage{hyperref}
\usepackage{cleveref}
\DeclareMathOperator{\Map}{Map}
\DeclareMathOperator{\Fun}{Fun}
\newcommand{\iso}{\xrightarrow{\sim}}
\newtheorem{lemma}{Supplementary Lemma}
\newtheorem{corollary}[lemma]{Supplementary Corollary}
\crefname{lemma}{Supplementary Lemma}{Supplementary Lemmas}
\Crefname{lemma}{Supplementary Lemma}{Supplementary Lemmas}
\crefname{corollary}{Supplementary Corollary}{Supplementary Corollaries}
\Crefname{corollary}{Supplementary Corollary}{Supplementary Corollaries}

\begin{document}
\section*{Fiber squares and diagonal squares}

The \href{pullbacks.html\#post:Representability_Of_Identifications}{representability axiom}
states a fiber square for the primitive identification anima.
The diagonal formulation is a consequence. Its proof compares entire cones,
including the identification expressing commutativity.

Fix absolute objects $a,b\colon\star\to E$, and put
$F=(a,b)\colon\star\to E\times E$ and
$\Delta=(\mathrm{id}_E,\mathrm{id}_E)\colon E\to E\times E$.
Vertical composition applies the rightmost factor first; $\star$ denotes whiskering.
All cones below lie in the absolute Chapter 1 setting.

\begin{lemma}
\label[lemma]{lem:Supplementary_Diagonal_Coordinates}
A cone $t=(u,v,\theta)$ over $(F,\Delta)$ on $T$ has two normalized
coordinates $x\colon au=v$ and $y\colon bu=v$, obtained by projecting
$\theta\colon Fu=\Delta v$ and inserting the product and unit comparisons.
It determines a cone $R(t)$ over $(a,b)$ with both legs $u$ and matching
$y^{-1}\cdot x$.
Cone comparisons induce comparisons of the recovered cones. Conversely,
a comparison $R(s)\to R(t)$ supplies a comparison $s\to t$.
\end{lemma}
\begin{proof}[Proof idea]
Project the matching equation of a cone comparison to its two coordinates.
The two resulting equations imply the equation for $y^{-1}\cdot x$.

Conversely, write the two leg comparisons of a recovered comparison as
$l,r\colon u_s=u_t$. Terminal uniqueness identifies $r$ with $l$.
Choose the comparison between the original right legs to be
\[
 y_t\cdot(b\star l)\cdot y_s^{-1}\colon v_s=v_t.
\]
Its second coordinate equation follows by cancellation. Its first coordinate
equation is the recovered matching equation, after replacing $r$ by $l$.
Detection by the two product projections gives the full cone comparison.
This establishes preservation and reflection of the existence of comparisons.
\end{proof}

\begin{lemma}
\label[lemma]{lem:Supplementary_Diagonal_Restriction}
For $r\colon S\to T$, recovery commutes with restriction by a cone
comparison $R(r^*t)\to r^*R(t)$ whose two leg comparisons are identities.
\end{lemma}
\begin{proof}[Proof idea]
With $\alpha_{r,u,a}\colon(au)r=a(ur)$, the normalized coordinates satisfy
\[
 x_{r^*t}=(x_t\star r)\cdot\alpha_{r,u,a}^{-1},\qquad
 y_{r^*t}=(y_t\star r)\cdot\alpha_{r,u,b}^{-1}.
\]
Insert these formulas into $y_{r^*t}^{-1}\cdot x_{r^*t}$.
The inverse-of-composite and whiskering comparisons yield
\[
 \alpha_{r,u,b}\cdot((y_t^{-1}\cdot x_t)\star r)
                       \cdot\alpha_{r,u,a}^{-1},
\]
which is the matching of the restricted fiber cone. Unit comparisons
turn this identification of matchings into the claimed cone comparison.
\end{proof}

\begin{lemma}
\label[lemma]{lem:Supplementary_Diagonal_Representability}
Let $I=(a=b)$ and let $\varepsilon\colon a!=b!$ be the specified
identification supplied by representability. There is a pullback cone over
$(\Delta,F)$ on $I$, with legs $a!\colon I\to E$ and $!\colon I\to\star$,
whose matching has normalized coordinates
\[
 (\mathrm{id}_{a!},\varepsilon).
\]
\end{lemma}
\begin{proof}[Proof idea]
First work with the reversed cospan $(F,\Delta)$. Choose its matching
with normalized coordinates $(\mathrm{id}_{a!},\varepsilon^{-1})$.
Recovery gives a fiber cone whose matching is
$(\varepsilon^{-1})^{-1}$. Representability, pullback symmetry, and
transport of the matching show that this recovered cone is a pullback.

Recover the chosen pullback cone over $(F,\Delta)$ and factor it through
this fiber cone. \Cref{lem:Supplementary_Diagonal_Restriction} compares
recovery before and after restriction, and
\Cref{lem:Supplementary_Diagonal_Coordinates} lifts the resulting cone
comparison back to the diagonal coordinates. Applying the same reflection
to comparisons of factorizations proves the pullback criterion.

Finally exchange the two legs. The matching is inverted; its normalized
coordinates become $(\mathrm{id}_{a!},\varepsilon)$ by the inverse and
double-inverse comparisons. Thus the result concerns this specified
matching, and introduces no further representability axiom.
\end{proof}
\end{document}
